\chapter{Generative Recursion}\label{ch:generative}

\begin{definition}[Generative recursion]
A recursive function uses \textbf{generative recursion} if at least one
recursive call is on a value that is \emph{generated} by computation, rather
than being a structural part of the input (like \code{(rest lst)} or
\code{(sub1 n)}).
\end{definition}

Generative recursion is more powerful but has two costs: the code no longer
follows from a template, and \textbf{termination} is no longer guaranteed.
The design recipe is modified: think of the algorithm first (how to generate
smaller subproblems and combine their solutions), and argue termination.

\section{Euclid's algorithm}
For $m, n \in \N$ with $n > 0$: $\gcd(m, 0) = m$ and
$\gcd(m, n) = \gcd(n, m \bmod n)$.

\begin{lstlisting}[style=racket]
;; (euclid-gcd m n) produces the greatest common divisor of m and n.
;; Examples:
(check-expect (euclid-gcd 24 18) 6)
(check-expect (euclid-gcd 7 0) 7)

;; euclid-gcd: Nat Nat -> Nat
(define (euclid-gcd m n)
  (cond [(zero? n) m]
        [else (euclid-gcd n (remainder m n))]))
\end{lstlisting}

\paragraph{Trace.}
\begin{lstlisting}[style=racket]
(euclid-gcd 24 18)
=> (euclid-gcd 18 6)      ; (remainder 24 18) = 6
=> (euclid-gcd 6 0)       ; (remainder 18 6) = 0
=> 6
\end{lstlisting}

\paragraph{Correctness.} Any common divisor of $m$ and $n$ divides
$m - qn = m \bmod n$, and conversely any common divisor of $n$ and $m \bmod n$
divides $m$. So $\gcd(m,n) = \gcd(n, m \bmod n)$.

\paragraph{Termination.} The second argument strictly decreases
($0 \le m \bmod n < n$) and is a natural number, so it reaches 0 after finitely
many steps. (In fact after at most about $2\log_2 n$ steps, since the second
argument at least halves every two calls.)

\begin{important}[Termination arguments]
To prove a generative recursion terminates, find a \textbf{measure}: a
natural number computed from the arguments that strictly decreases at every
recursive call. Since naturals cannot decrease forever, recursion must stop.
\end{important}

\subsection{A function that may not terminate}
The Collatz function:
\begin{lstlisting}[style=racket]
;; collatz: Nat -> Nat     requires: n >= 1
(define (collatz n)
  (cond [(= n 1) 1]
        [(even? n) (collatz (/ n 2))]
        [else (collatz (+ 1 (* 3 n)))]))
\end{lstlisting}
No one knows whether this terminates for every $n \ge 1$ --- it is a famous
open problem. This illustrates why termination needs an argument.

\section{Quicksort}
Quicksort (C.~A.~R.~Hoare, 1960) is a divide-and-conquer sorting algorithm:
\begin{enumerate}
  \item If the list is empty, it is sorted.
  \item Otherwise pick a \textbf{pivot} (the first element).
  \item Partition the rest into elements less than the pivot and the others.
  \item Sort each part recursively and append: smaller, pivot, larger.
\end{enumerate}

\begin{lstlisting}[style=racket]
;; (quick-sort lon) produces lon sorted in non-decreasing order.
;; Examples:
(check-expect (quick-sort empty) empty)
(check-expect (quick-sort '(3 1 4 1 5 9 2 6)) '(1 1 2 3 4 5 6 9))

;; quick-sort: (listof Num) -> (listof Num)
(define (quick-sort lon)
  (cond [(empty? lon) empty]
        [else
         (local [(define pivot (first lon))
                 (define less (filter (lambda (x) (< x pivot)) (rest lon)))
                 (define greater (filter (lambda (x) (>= x pivot)) (rest lon)))]
           (append (quick-sort less)
                   (list pivot)
                   (quick-sort greater)))]))
\end{lstlisting}

The recursive calls are on \code{less} and \code{greater}, lists computed by
\code{filter} --- this is generative recursion.

\paragraph{Condensed trace.}
\begin{lstlisting}[style=racket]
(quick-sort (list 3 5 1 4))
=> (append (quick-sort (list 1)) (list 3) (quick-sort (list 5 4)))
=> (append (append (quick-sort empty) (list 1) (quick-sort empty))
           (list 3) (quick-sort (list 5 4)))
=> (append (list 1) (list 3) (quick-sort (list 5 4)))
=> (append (list 1) (list 3)
           (append (quick-sort (list 4)) (list 5) (quick-sort empty)))
=> (append (list 1) (list 3) (append (list 4) (list 5) empty))
=> (list 1 3 4 5)
\end{lstlisting}

\paragraph{Termination.} Measure: length of the list. Both \code{less} and
\code{greater} are sublists of \code{(rest lon)}, so each has length strictly
less than \code{(length lon)}.

\paragraph{Efficiency.} Partitioning costs time proportional to the list length.
If the pivot splits the list roughly in half each time, there are about
$\log_2 n$ levels, each with $\Theta(n)$ total work: $\Theta(n \log n)$. If the
list is already sorted, \code{less} is always empty and \code{greater} has
length $n-1$: $n + (n-1) + \cdots = \Theta(n^2)$. A common bug is putting
elements equal to the pivot in neither part (losing duplicates) or in a part
with the pivot itself (non-termination).

\section{Mergesort}
Mergesort splits the list in half, sorts each half recursively, and uses
\code{merge} from Chapter 8. It is always $\Theta(n \log n)$.
\begin{lstlisting}[style=racket]
;; merge-sort: (listof Num) -> (listof Num)
(define (merge-sort lon)
  (cond [(or (empty? lon) (empty? (rest lon))) lon]
        [else
         (local [(define half (quotient (length lon) 2))]
           (merge (merge-sort (take lon half))
                  (merge-sort (drop lon half))))]))

;; (take lst n) produces the first n elements; (drop lst n) the rest.
;; take: (listof X) Nat -> (listof X)
(define (take lst n)
  (cond [(zero? n) empty]
        [else (cons (first lst) (take (rest lst) (sub1 n)))]))
;; drop: (listof X) Nat -> (listof X)
(define (drop lst n)
  (cond [(zero? n) lst]
        [else (drop (rest lst) (sub1 n))]))
\end{lstlisting}
Termination: for lists of length $\ge 2$, both halves are strictly shorter.

\section{Other examples of generative recursion}

\subsection{Binary search for a root (bisection)}
To find $x$ with $f(x) \approx 0$ on $[lo, hi]$ where $f(lo)$ and $f(hi)$
have opposite signs, repeatedly halve the interval:
\begin{lstlisting}[style=racket]
;; (find-root f lo hi tol) produces an approximate root of f in [lo, hi].
;; find-root: (Num -> Num) Num Num Num -> Num
;;   requires: (f lo) and (f hi) have opposite signs; tol > 0
(define (find-root f lo hi tol)
  (local [(define mid (/ (+ lo hi) 2))]
    (cond [(< (- hi lo) tol) mid]
          [(<= (* (f lo) (f mid)) 0) (find-root f lo mid tol)]
          [else (find-root f mid hi tol)])))

(check-within (find-root (lambda (x) (- (* x x) 2)) 0 2 0.0001)
              (sqrt 2) 0.001)
\end{lstlisting}
Termination: the interval width halves each call, so after
$\ceil{\log_2((hi-lo)/tol)}$ calls it is below the tolerance.

\subsection{Processing strings in chunks}
Splitting a list of characters into words, where each recursive call is on
the characters after the next space, is generative: the amount removed depends
on the data.

\section{Design recipe for generative recursion}
\begin{enumerate}
  \item Identify trivially solvable cases and their solutions.
  \item Describe how to generate one or more smaller problems.
  \item Describe how to combine their solutions.
  \item Give a termination argument (a decreasing measure) in a comment.
  \item Test the trivial cases and cases that exercise each kind of split; include inputs that might cause non-termination (duplicates, sorted input).
\end{enumerate}

\section{Summary}
\begin{itemize}
  \item Generative recursion recurses on computed values; it is not driven by a template.
  \item Termination must be argued with a decreasing natural-number measure.
  \item Euclid's gcd, quicksort, mergesort and bisection are the canonical examples.
  \item Quicksort is $\Theta(n \log n)$ on average and $\Theta(n^2)$ in the worst case; mergesort is always $\Theta(n \log n)$.
\end{itemize}
